\chapter{IMPLEMENTATION}

\label{ch:implementation}

This chapter describes the artifact as it stands, increment by increment. Each section states the
property being protected, the mechanism that protects it, and the boundary the mechanism does not
cross. Decisions reversed during development are recorded at the end, because the reversals are
part of the evidence: the first design was a single graph view with browser-local state, and almost
none of that layer survives.

\section{Accounts, Sessions and Scope}

Authentication issues a JSON Web Token into an {\tt httpOnly}, {\tt SameSite=Lax} cookie marked
{\tt Secure} in production. The client never reads the token, so a script injected into the page
cannot exfiltrate it. Because a token is a snapshot, the server does not trust what it carries: on
every request the middleware re-reads the account from the database and attaches the current role
and program context. A token issued while an account was a student cannot act as an administrator
after a role change.

The second half of the mechanism is the expected scope. A mutation such as completing a course,
recording a grade, casting a rating or saving a semester may carry an optional structure naming the
owner and the program context the client observed. The server compares that expectation with the
values it reads \emph{inside the transaction that performs the write}. Equality proceeds; a
mismatch is refused with a conflict and nothing is written. An expected scope is never persisted.

Two consequences are worth stating. A stale tab cannot silently write into a scope the user no
longer occupies. And the mechanism is deliberately backward-compatible: a request that omits the
expectation is treated as a legacy client, so the client could adopt the precondition one screen at
a time rather than in a single flag day.

The client half of the same problem is the session-expiry race. A protected read refused with 401
once ended the session even when the user had already signed in again while that read was in
flight. Each request now records the session generation it was sent under; a refusal belonging to
an earlier generation is ignored, and a refusal from the current generation still ends the session.

\section{Program Context and the Reference Import}

The program model separates four ideas a single table would conflate. A \emph{course} is a global
identity with a code, a name and a credit count. A \emph{membership} says the course belongs to a
program and is unique per pair. A \emph{placement} says where the program lists the course --
academic year, semester, elective group, selection count -- and is not unique, because the same
course may be listed twice with different source order. A \emph{prerequisite edge} connects two
memberships of the same program.

Keeping courses global and placement local is what makes multiple programs possible at all. A shared
mathematics course exists once, so its rating average serves every program, while its year and
semester differ per program and therefore live in placements.

The import runs in three independently inspectable stages: a verifier that proves the source is
losslessly representable, a preflight that inspects the target database, and an idempotent backfill.
Two source irregularities are modelled rather than normalised away. A blank course code appears in
the source as an explicit free-elective requirement, so it is stored as a requirement with its own
credits, not as a course. A course listed in several elective groups becomes several placements
sharing one membership, so a completed course is claimed by exactly one group while remaining
hidden in the others.

The program records its provenance and admits its uncertainty. The stored credit total is nullable
and null in the reference import, so readers resolve an unknown percentage rather than dividing by
an assumed target. The final-year branch -- a thesis track above a grade threshold and a coursework
track below it -- is retained as a declared reference rule, and the interface says so.

An administrator can attach or clear a reference program for a student account, which is what makes
the context-aware screens demonstrable. The control is deliberately narrow: it writes one column in
one transaction, refuses a stale precondition, returns a {\tt changed} flag so an identical retry
recovers instead of conflicting, writes no record, grade, rating or plan, and labels every reply as
a reference attachment rather than a cohort assignment.

\section{Progress, Prerequisites and Cascades}

Progress is one row per account and course, with a status and, where it applies, the elective group
that claimed the course. The client keeps an optimistic view so the interface responds immediately,
but the stored truth is the server's.

Completion is validated on the server against the program's mandatory prerequisite edges. Every
declared edge is mandatory: a legacy corequisite or strictness column is preserved as provenance
but does not soften the rule, which removed an earlier client array of exceptions that duplicated,
and disagreed with, the database.

Uncompletion cascades, because leaving a dependent completed would violate the system's own rule. On
the server the closure is computed inside the same serializable transaction as the status change, so
the invariant holds at commit rather than eventually. On the client the same closure is computed
optimistically from a reverse dependency map built over the loaded edges,
\[
  \mathit{dependencyMap}[p] = \{\, c \mid (c, p) \in \mathit{prerequisites} \,\},
\]
so the interface updates before the reply arrives. The client result is a prediction; if the
server's returned snapshot differs, the snapshot wins.

\begin{algorithm}[H]
\caption{Optimistic dependent removal on the client}
\KwIn{reverse map $D$, completed record $C$, course $c$}
\KwOut{updated completed record $C'$}
$queue \leftarrow [c]$\;
\While{$queue$ is not empty}{
  $current \leftarrow \text{pop}(queue)$\;
  $C \leftarrow C \setminus \{current\}$\;
  \ForEach{$dependent \in D[current]$}{
    \If{$dependent \in C$}{
      $queue.\text{push}(dependent)$\;
    }
  }
}
\Return $C$\;
\end{algorithm}

The same closure is applied when a completed course is converted to planned: the course leaves the
completed set, enters the planned set, and its completed dependents are removed with it.

Archive import follows the same discipline. An account's older browser-local selections can be
imported, but the import validates membership, prerequisites and cycles atomically and refuses a
selection the assigned program does not permit, rather than accepting a client's opinion about what
the student has done.

\section{The Unassigned Curriculum Map}

An account with no program context sees the editable curriculum map. It is a semester-column layout
that mirrors how the program is published: eight standard semesters plus a summer column where the
catalogue declares one. Each course is a card showing its code, name, credits and requirement
category, and its state is carried by more than colour: completed, planned, locked because a
mandatory prerequisite is unmet, recommended for the current window, highlighted as a prerequisite
of a hovered locked course, or dimmed as a duplicate occurrence of a course claimed elsewhere.

A left click toggles completion and a right click toggles planning; both actions are also available
from the keyboard, and every control names what it does in text a screen reader reads rather than
only in a tooltip. Completing an elective-group course records the group that claimed it, so the
course is counted once and hidden from the other groups that also list it:
\[
  \mathit{completedIds}[\mathit{courseId}] = \mathit{electiveGroupName}.
\]

Clicking a group label opens a sidebar listing that group's members with their own states. The map
rescales when the sidebar opens or closes and deliberately does not rescale when the student
switches between groups while the sidebar is already open, because that transition does not change
the available width.

The intensity control sets the credit band used for the recommended set: nine, fifteen, eighteen or
twenty-one credits per semester for the low, normal, high and maximum settings. The recommended set
is computed from the loaded catalogue and the current completion state. The final-year branch
determines whether the thesis course or the alternative coursework is offered, and the reference
target used for the progress figure follows the same branch, with physical training excluded from
earned credits but included in planning. Where the target is unknown the percentage is withheld
rather than assumed.

\section{Grades, Averages and Per-Semester Records}

A grade is an immutable attempt, not a field on a record. Each attempt stores a course, a score from
0 to 100, optional term metadata, and a request identifier unique per account. A retake appends; it
never overwrites. An identical retry returns the original attempt, so a lost reply is safe to
repeat, and a different payload under a used identifier is refused rather than applied.

The cumulative average is credit-weighted over the highest score recorded for each course, with
physical training excluded because it carries no grade weight in the program. Arithmetic is decimal,
not binary floating point, because the threshold that selects the thesis track is an exact
comparison and a value that should sit exactly at the threshold must not fall below it through
rounding.

Per-semester averages are a separate, deliberately narrower claim. A term exists only when an
attempt stores a year and a semester; creation time and a saved plan's placement are never used to
infer one. Each term keeps its own highest score per course, so a retake can count in more than one
term, and the weighting reuses the cumulative evaluator. Attempts without dates are counted and
reported as undated rather than assigned to a guessed term. The interface states that these labels
are self-recorded reference metadata, not calendar or registration evidence.

\section{Difficulty Ratings and Shrinkage}

Students may rate a completed course from one to five. Write access is gated on a completion record
in the same snapshot, limited by a durable hourly quota per account, and idempotent by vote: an
unchanged repeat is acknowledged rather than counted twice. The cached average and count on the
course row are updated in the same transaction as the vote.

The displayed difficulty is a shrinkage estimate. With $v$ the course's average vote, $n$ its vote
count, $m$ the program-level prior mean and $C$ a prior strength expressed as pseudo-votes,
\[
  \hat{d} \;=\; \frac{v n + m C}{n + C}.
\]
At $n = 0$ the estimate is exactly the prior; as $n$ grows the course's own evidence dominates. When
a program has no rated courses the prior falls back to the seed difficulty mean of its members, and
the interface distinguishes the two sources. The vote count is shown next to every estimate, and a
course with no votes is described as having no votes rather than as being of average difficulty.

\section{Workload Analysis and Recommendations}

Workload analysis and semester recommendations share one resolver. It reads the account's program,
placements, prerequisite edges, progress, highest numeric scores and rating prior in a single
snapshot, and rejects a selection that is not a member of that program instead of falling back to
the global catalogue. Recommendations rank eligible courses by the shrunk difficulty estimate and a
bounded numeric-grade fit, then budget credits by the selected intensity. Courses that cannot be
scheduled -- because a mandatory prerequisite is unmet, or because the term label contradicts the
placement -- are reported with their reason instead of being silently dropped, and the interface
states that no graduation estimate is available while the program's own rules remain unverified.

The final-year branch is applied before ranking, not after: a student on the thesis path is not
offered the coursework the alternative path would require, and the reverse. Where eligibility is
unknown -- a program without a numeric average, or a guest -- both options are kept and the
interface says so rather than choosing one.

\section{Saved Plans and the Reference-Selection Editor}

Saved plans are browsed in three steps: plan metadata, then its stored semesters, then the courses
of one semester. Each read is bounded and keyset paginated, and each returns enough of its scope
that the client can verify the reply belongs to the object it asked for.

Editing a saved semester is the most constrained operation in the system and the one with the most
recovery. The write must confirm, inside one serializable transaction, that the caller owns the
plan, that the plan's program is still the caller's current program, that the semester's stored
revision still matches what the client saw, that every selected course is a current member with a
placement, that every mandatory prerequisite is completed, and what the caller's numeric average is
if the selection depends on the branch. Credits and difficulty totals are computed from that same
snapshot, and the revision increment and the immutable receipt commit with the selection. An exact
retry of a committed request is acknowledged from the receipt before the current-context checks, so
a lost reply does not become a lost edit; a changed payload under the same key conflicts.

Around that transaction the client provides three more layers. An unsent draft is stored per tab,
bound to the exact account, program, plan, semester and base revision, and restores only against a
candidate that still matches; a newer saved revision makes the draft discard-only rather than
rebased. Proposed courses can be checked before saving, reporting per course which prerequisite is
missing, whether the placement is wrong, whether the numeric average is unknown, and whether the
course is placed at all -- and the interface states that this diagnostic never authorises the later
save, which rechecks everything. The stored list order can be changed with explicit move controls,
and the interface states that order does not alter academic chronology, prerequisite satisfaction or
allocation priority.

\section{Administration: Resources, Demand and Capacity}

The administrative side is explicitly a simulation and is labelled as such everywhere.

Resource settings are scoped to a program, a semester and a year, and consist of a professor count,
a classroom count, a laboratory count, a maximum section size, and optional per-course overrides.
Each save carries the revision the administrator observed; the update is a compare-and-set, so a
concurrent save produces a conflict rather than a lost update. Absent settings are unknown, an
explicit zero is zero, and every diagnostic keeps the two apart.

Demand is counted from the current assigned cohort's planned selections, with each student counted
once per course. A planned selection that is not a member of the program is excluded and the
exclusion is reported. Capacity comparisons use declared seats where an override exists and report
unknown where it does not. A classroom-seat proxy is presented separately and never copied into
every course, because one room can host only one simultaneous section.

\section{Allocation Models and the Explicit Executor}

Two allocation models exist and are versioned contracts rather than code paths. The earlier model
allocates one course per student in a single scarcity-ordered round. The later model allocates
several courses per student across rounds against one shared section-and-seat ledger, with a
per-student credit target and explicit stop reasons. Both order students by how constrained they
are, break ties deterministically, and report individual outcomes separately from the aggregate.

Each capture stores the model identifier, the captured policies, the aggregate result and the
private per-student results, and each read replays the stored contract rather than re-deriving it
from current configuration. That is what makes history meaningful: a capture taken under one
resource policy remains readable after the policy changes, and a corrupt or mismatched record fails
closed instead of being partially displayed.

Execution is explicit. A queued request names its scenario and its actor; one job is executed either
by a command-line invocation or by a protected administrator action, inside one transaction that
claims the job, verifies the caller, captures the result and writes the terminal outcome together.
There is no daemon, no queue drain and no timer. A crash leaves the request pending rather than half
executed, and an exact retry either recovers the original outcome or finds the job still pending.

\section{Client Interface and Accessibility}

The interface is a set of routes -- curriculum, grades, ratings, planner, administration -- behind
an authentication guard, with an administrator guard on the administrative surface.

The curriculum route branches on context. With no program the editable map is shown; with an
assigned program the screen becomes a read-only reference: placements, mandatory prerequisites,
difficulty estimates with real vote counts, an explicit list of unresolved requirements, and a
progress summary that separates current member data from preserved history. The branch is
deliberate. Editing inside an assigned program needs a context-keyed cache and a claim-transfer
policy that were not built, so the interface refuses to pretend otherwise.

The planner follows the same shape, with the legacy analysis for guests and the context preview for
assigned accounts. The grades screen shows the average, the retake history, a highest-score
distribution and, where the data supports it, the per-semester table. The ratings screen lists only
courses the account has completed and keeps a submitted vote's retry identity until the server
confirms it.

Because the interface is dense, it was audited rather than assumed correct. The audit found a map
that faded entirely when a locked course was hovered, an elective panel more than forty tab stops
from the control that opened it and indifferent to Escape, a save that dropped keyboard focus to the
document body because the focused control became disabled, zoom buttons that looked available at
their limits, a crash screen with no exit, an error message that left focus at the top of the
document, a rejected score that did not point at the field that needed fixing, and the session-
expiry race described above.

\section{Reversed Decisions}

Three early decisions were reversed on evidence, and the reversals explain why the artifact looks
the way it does.

\textbf{The graph view was deleted.} An interactive node-and-edge curriculum graph was built on a
third-party layout library and then removed: it duplicated the semester-column layout, was the only
consumer of two dependencies, and could not represent repeated elective placements without inventing
positions. The semester columns, which mirror how the program is published, replaced it.

\textbf{Browser-local state was replaced by server state.} The first version persisted completions
and plans in browser storage, which made progress device-specific and unverifiable. Browser storage
now remains only as a confirmed cache for an account with no program context, and is refused as
truth once a program is attached.

\textbf{A duplicated rule array was removed.} The planner originally carried its own list of
semester ordering rules alongside the database prerequisites, and the two disagreed. The rules now
live only in the database, and every edge is mandatory.
